\chapter{Kiểm thử và đánh giá hệ thống}

\section{Giới thiệu}

Sau khi hoàn thành quá trình cài đặt và triển khai, hệ thống AITasker được tiến
hành kiểm thử nhằm đánh giá mức độ đáp ứng các yêu cầu chức năng và phi chức
năng đã được đặc tả. Quá trình kiểm thử giúp phát hiện các lỗi còn tồn tại,
đánh giá tính ổn định, hiệu năng và khả năng sử dụng của hệ thống trước khi đưa
vào sử dụng.

Trong phạm vi đồ án, nhóm tập trung thực hiện kiểm thử chức năng (Functional
Testing) kết hợp với kiểm thử tích hợp (Integration Testing) và kiểm thử giao
diện người dùng.

%============================================================

\section{Môi trường kiểm thử}

Quá trình kiểm thử được thực hiện trong môi trường trình bày ở
Bảng~\ref{tab:test_environment}.

\begin{table}[H]
\centering
\caption{Môi trường kiểm thử}
\label{tab:test_environment}

\begin{tabular}{|p{6cm}|p{8cm}|}
\hline
\textbf{Thành phần} & \textbf{Thông tin} \\ \hline

Hệ điều hành & Windows 11 \\ \hline

Backend & FastAPI \\ \hline

Frontend & ReactJS \\ \hline

Database & PostgreSQL \\ \hline

API Testing & Swagger UI, Postman \\ \hline

IDE & Visual Studio Code \\ \hline

Trình duyệt & Google Chrome \\ \hline

\end{tabular}

\end{table}

%============================================================

\section{Phương pháp kiểm thử}

Các phương pháp kiểm thử được áp dụng bao gồm:

\begin{itemize}

\item Kiểm thử chức năng.

\item Kiểm thử giao diện.

\item Kiểm thử tích hợp.

\item Kiểm thử API.

\item Kiểm thử phân quyền.

\item Kiểm thử dữ liệu đầu vào.

\end{itemize}

%============================================================

\section{Các chức năng được kiểm thử}

Nhóm tiến hành kiểm thử các chức năng chính của hệ thống:

\begin{enumerate}

\item Đăng ký tài khoản.

\item Đăng nhập.

\item Quản lý hồ sơ.

\item Tạo Project.

\item Cập nhật Project.

\item Gửi Proposal.

\item Chấp nhận Proposal.

\item Tạo Contract.

\item Quản lý Milestone.

\item Thanh toán.

\item Đánh giá.

\item Dashboard.

\item Notification.

\end{enumerate}

%============================================================

\section{Các trường hợp kiểm thử}

\subsection{Kiểm thử chức năng đăng ký}

\begin{table}[H]
\centering
\caption{Kiểm thử chức năng đăng ký}
\begin{tabular}{|p{2cm}|p{4cm}|p{4cm}|p{4cm}|}
\hline

\textbf{Test ID}
&
\textbf{Dữ liệu đầu vào}
&
\textbf{Kết quả mong đợi}
&
\textbf{Kết quả}
\\
\hline

TC01
&
Thông tin hợp lệ
&
Đăng ký thành công
&
Đạt
\\
\hline

TC02
&
Email đã tồn tại
&
Thông báo lỗi
&
Đạt
\\
\hline

TC03
&
Mật khẩu dưới 8 ký tự
&
Thông báo lỗi
&
Đạt
\\
\hline

\end{tabular}

\end{table}

%============================================================

\subsection{Kiểm thử chức năng đăng nhập}

\begin{table}[H]
\centering
\caption{Kiểm thử chức năng đăng nhập}

\begin{tabular}{|p{2cm}|p{4cm}|p{4cm}|p{4cm}|}

\hline

TC04
&
Email và mật khẩu đúng
&
Đăng nhập thành công
&
Đạt
\\
\hline

TC05
&
Sai mật khẩu
&
Thông báo lỗi
&
Đạt
\\
\hline

TC06
&
Email không tồn tại
&
Thông báo lỗi
&
Đạt
\\
\hline

\end{tabular}

\end{table}

%============================================================

\subsection{Kiểm thử chức năng tạo Project}

\begin{table}[H]

\centering

\caption{Kiểm thử chức năng tạo Project}

\begin{tabular}{|p{2cm}|p{4cm}|p{4cm}|p{4cm}|}

\hline

TC07
&
Thông tin hợp lệ
&
Project được tạo
&
Đạt
\\
\hline

TC08
&
Thiếu tiêu đề
&
Thông báo lỗi
&
Đạt
\\
\hline

TC09
&
Ngân sách âm
&
Thông báo lỗi
&
Đạt
\\
\hline

\end{tabular}

\end{table}

%============================================================

\subsection{Kiểm thử chức năng gửi Proposal}

\begin{table}[H]

\centering

\caption{Kiểm thử chức năng gửi Proposal}

\begin{tabular}{|p{2cm}|p{4cm}|p{4cm}|p{4cm}|}

\hline

TC10
&
Thông tin hợp lệ
&
Proposal được tạo
&
Đạt
\\
\hline

TC11
&
Đã gửi Proposal
&
Thông báo lỗi
&
Đạt
\\
\hline

\end{tabular}

\end{table}

%============================================================

\subsection{Kiểm thử chức năng thanh toán}

\begin{table}[H]

\centering

\caption{Kiểm thử chức năng thanh toán}

\begin{tabular}{|p{2cm}|p{4cm}|p{4cm}|p{4cm}|}

\hline

TC12
&
Milestone hợp lệ
&
Thanh toán thành công
&
Đạt
\\
\hline

TC13
&
Milestone chưa hoàn thành
&
Thông báo lỗi
&
Đạt
\\
\hline

\end{tabular}

\end{table}

%============================================================

\section{Kiểm thử API}

Nhóm sử dụng Swagger UI và Postman để kiểm thử các RESTful API.

Các API đã được kiểm thử gồm:

\begin{itemize}

\item Authentication API

\item User API

\item Enterprise API

\item Expert API

\item Project API

\item Proposal API

\item Contract API

\item Milestone API

\item Review API

\item Notification API

\item Dashboard API

\end{itemize}

Tất cả API đều trả về đúng định dạng JSON và mã trạng thái HTTP tương ứng.

%============================================================

\section{Kiểm thử giao diện}

Giao diện hệ thống được kiểm thử trên các trình duyệt:

\begin{itemize}

\item Google Chrome

\item Microsoft Edge

\item Mozilla Firefox

\end{itemize}

Kết quả cho thấy giao diện hiển thị đúng, bố cục hợp lý và hỗ trợ Responsive
trên các kích thước màn hình phổ biến.

%============================================================

\section{Kiểm thử phân quyền}

Các vai trò được kiểm thử gồm:

\begin{itemize}

\item Administrator

\item Enterprise

\item Expert

\end{itemize}

Kết quả kiểm thử cho thấy:

\begin{itemize}

\item Administrator có toàn quyền quản trị.

\item Enterprise chỉ được thao tác trên Project và Contract của mình.

\item Expert chỉ được gửi Proposal và thực hiện Contract.

\item Người dùng chưa đăng nhập không thể truy cập các API bảo vệ.

\end{itemize}

%============================================================

\section{Đánh giá hệ thống}

Sau quá trình kiểm thử, hệ thống đạt được các kết quả sau:

\begin{itemize}

\item Các chức năng chính hoạt động đúng yêu cầu.

\item Phân quyền hoạt động chính xác.

\item API phản hồi đúng dữ liệu.

\item JWT xác thực ổn định.

\item PostgreSQL lưu trữ dữ liệu chính xác.

\item Hệ thống vận hành ổn định.

\item Giao diện thân thiện.

\item Kiến trúc dễ mở rộng.

\end{itemize}

%============================================================

\section{Hạn chế}

Mặc dù hệ thống đã đáp ứng hầu hết các yêu cầu đặt ra, vẫn còn một số hạn chế:

\begin{itemize}

\item Chưa tích hợp cổng thanh toán thực tế.

\item Chưa triển khai AI Recommendation.

\item Chưa triển khai WebSocket cho thông báo thời gian thực.

\item Chưa có ứng dụng di động.

\item Chưa thực hiện kiểm thử tải (Load Testing) và kiểm thử hiệu năng chuyên sâu.

\end{itemize}

%============================================================

\section{Hướng phát triển}

Trong tương lai, hệ thống có thể được mở rộng theo các hướng sau:

\begin{itemize}

\item Tích hợp Stripe hoặc PayPal.

\item Tích hợp AI Recommendation System.

\item Tích hợp AI Chatbot nâng cao.

\item Triển khai Docker và Kubernetes.

\item Triển khai CI/CD.

\item Xây dựng ứng dụng Android và iOS.

\item Bổ sung hệ thống giám sát và cảnh báo.

\item Tối ưu hiệu năng cơ sở dữ liệu.

\end{itemize}

%============================================================

\section{Kết luận}

Qua quá trình kiểm thử, hệ thống AITasker đáp ứng đầy đủ các yêu cầu chức năng
và phần lớn các yêu cầu phi chức năng đã đề ra. Các chức năng chính hoạt động
ổn định, dữ liệu được quản lý chính xác và cơ chế phân quyền đảm bảo an toàn.
Mặc dù vẫn còn một số hạn chế cần tiếp tục cải thiện, hệ thống đã sẵn sàng để
phát triển và triển khai trong các giai đoạn tiếp theo.